\section{Lựa chọn mô hình triển khai}

Model triển khai được chọn từ các cấu hình không bị trội theo tập tiêu chí chất
lượng và hiệu năng. Trước hết, loại các cấu hình không đạt ngưỡng accuracy,
memory hoặc latency. Trong nhóm còn lại, ưu tiên cấu hình ổn định trên cả hai
thiết bị và không phụ thuộc vào một tối ưu khó tái lập.

Kết quả lựa chọn cần ghi rõ model version, pruning ratio, QAT configuration,
precision, input shape, checksum, engine tương ứng với từng device class và bộ
chỉ số benchmark. Model được chọn là đầu vào của Model Registry và kịch bản
Deployment trong các chương tiếp theo, không phải một tệp tách rời khỏi bằng
chứng thực nghiệm.

\begin{table}[H]
\centering
\caption{Quyết định lựa chọn model triển khai}
\label{tab:selected-model}
\begin{tabularx}{\textwidth}{L{4cm}Y}
\toprule
\textbf{Thuộc tính} & \textbf{Giá trị}\\
\midrule
Model version & \cbs{phiên bản được chọn}\\
Phương pháp tối ưu & \cbs{baseline/pruning/QAT/pruning + QAT}\\
Precision & \cbs{FP16/INT8}\\
Accuracy & \cbs{Precision, Recall, mAP}\\
Jetson Nano & \cbs{FPS, latency, memory, power}\\
Jetson Orin Nano & \cbs{FPS, latency, memory, power}\\
Lý do lựa chọn & \cbs{phân tích trade-off}\\
\bottomrule
\end{tabularx}
\end{table}